% Generated draft: requires mathematical and source review.
\documentclass[11pt]{article}
\usepackage[margin=25mm]{geometry}
\usepackage{fontspec}
\setmainfont{Latin Modern Roman}
\usepackage{amsmath,amssymb,mathtools,unicode-math}
\setmathfont{Latin Modern Math}
\usepackage{enumitem,needspace,xurl,hyperref}
\setlength{\parindent}{0pt}
\setlength{\parskip}{5pt}
\setlength{\emergencystretch}{3em}
\newcommand{\N}{\mathbb{N}}
\newcommand{\Z}{\mathbb{Z}}
\newcommand{\Q}{\mathbb{Q}}
\newcommand{\R}{\mathbb{R}}
\newcommand{\C}{\mathbb{C}}
\newcommand{\F}{\mathbb{F}}
\newcommand{\A}{\mathbb{A}}
\newcommand{\PP}{\mathbb{P}}
\newcommand{\E}{\mathbb{E}}
\newcommand{\eps}{\varepsilon}
\newcommand{\dd}{\,\mathrm{d}}
\begin{document}
\section*{UnsolvedMath ID: 700065}
\textbf{Source problem code:} catalogue entry
\subsection{Definitions and mathematical statement}
A topological group is a group \(G\) equipped with a topology \(\tau\) for which multiplication and inversion are continuous. A subgroup \(H\leq G\) is given the subspace topology \(\tau|_H\); it is discrete if every singleton is open in that topology.

The supplied statement asks to describe all non-discrete topological groups whose closed subgroups are discrete. Literally, its defining condition is
\[
\tau\neq\mathcal P(G),\qquad
\forall H\leq G,\quad
\bigl(H\text{ is closed in }G\bigr)
\Longrightarrow
\bigl(\tau|_H=\mathcal P(H)\bigr),
\]
where \(\mathcal P(X)\) denotes the power set of \(X\). This formulation has an unresolved incompatibility: \(G\) itself is a closed subgroup, so the displayed universal condition requires \(G\) to be discrete, contrary to the first requirement.

A possible intended question is to describe all non-discrete topological groups satisfying
\[
\forall H<G,\quad
\bigl(H\text{ is closed in }G\bigr)
\Longrightarrow
\bigl(H\text{ is discrete in its subspace topology}\bigr),
\]
where \(H<G\) means a proper subgroup. The supplied record does not include the word ``proper,'' so this is an explicitly provisional interpretation requiring source confirmation. The source calls the intended groups ``minimal.'' It does not specify a separation axiom, and in particular does not explicitly require the topology to be Hausdorff.

The source further states that the locally compact case can be described without much effort and suggests that the general case is probably complicated. These are supplied contextual assessments; neither a classification nor supporting hypotheses for the locally compact assertion are provided.
\subsection{Short English statement}
Platonov asks for a description of certain non-discrete topological groups. The wording requires every closed subgroup to be discrete, but includes the whole group and is therefore incompatible as written. Restricting the condition to proper closed subgroups is a possible intended reading requiring confirmation. The separation assumptions also remain unspecified.
\subsection{Sources}
\begin{itemize}
\item Supplied local UnsolvedMath registry record: permanent ID \(700065\), exact source problem code \texttt{catalogue entry}, title ``Kourovka Notebook Problem 1.74,'' upstream revision \texttt{d371ee917ac9647fbb06cd6fc529deec6f3dda82}. The record labels the problem ``open''; this is supplied registry status, not an independently verified mathematical assessment.
\item Attributed to V. P. Platonov, proposed in 1965. The supplied background identifies \emph{The Kourovka Notebook (Unsolved Problems in Group Theory)}, original problem 1.74, first issue (1965), as reproduced in the 21st-issue edition (2026), \texttt{arXiv:1401.0300v48}.
\item The supplied collection description names E. I. Khukhro and V. D. Mazurov as editors and describes coverage of Issues 1--21 (1965--2026). The cited edition and its wording have not been independently inspected for this draft.
\end{itemize}
\end{document}
